%% GENERATED FILE -- do not hand-edit.
%% SOURCE: experiments/database/scripts/build_candidate_datasets_table.py
\begin{landscape}
\begingroup
\footnotesize
\setlength{\tabcolsep}{3pt}
\renewcommand{\arraystretch}{1.15}
\begin{longtable}{@{}r@{\hspace{3pt}} L{40mm} L{11mm} L{20mm} L{18mm} r@{\hspace{4pt}} r@{\hspace{4pt}} L{19mm} L{14mm} L{81mm}@{}}
\caption[]{The final fifty, then ten extra in rank order in case a row above proves
unreachable. Every stat the ranking used is here beside the reason. \emph{Band, tier}:
\emph{P} is the perturb-seq band, \emph{M} molecular layers, \emph{S} scale, applied before
tier (Sec.~\ref{sec:rule}); within a band the tier rule orders, then measurements.
\emph{Genotypes} and \emph{Env} are the perturbation and condition axes. \emph{Inst.} is
genotype$\times$environment records, $\dagger$ where it is the product of the two axes
rather than a reported count and $\ddagger$ where it is an order-of-magnitude estimate.
\emph{Meas.} is instances times phenotype dimensionality, the quantity rows are ranked on.
\emph{Sequence basis} is the route to each strain's total genomic content; a row with no
route is excluded (Table~\ref{tab:excluded}). \emph{Time} names the row's time dimension
where it has one; a dash means steady state or endpoint. A $\bullet$ marks a row high on a
Perturb-seq axis; superscript \textbf{B} marks a row blocked on data access, \textbf{L} one
with a loader in flight. Citations, links and data locations are in
Table~\ref{tab:sources}; joins in Table~\ref{tab:synergies}.}
\label{tab:final}\\
\toprule
\textbf{\#} & \textbf{Dataset (class; phenotype)} & \textbf{Band, tier} & \textbf{Genotypes} & \textbf{Env} & \textbf{Inst.} & \textbf{Meas.} & \textbf{Sequence basis} & \textbf{Time} & \textbf{Why} \\
\midrule
\endfirsthead
\multicolumn{10}{@{}l}{\footnotesize\emph{Table~\ref{tab:final}, continued}}\\
\toprule
\textbf{\#} & \textbf{Dataset (class; phenotype)} & \textbf{Band, tier} & \textbf{Genotypes} & \textbf{Env} & \textbf{Inst.} & \textbf{Meas.} & \textbf{Sequence basis} & \textbf{Time} & \textbf{Why} \\
\midrule
\endhead
\bottomrule
\endfoot
1 & \textbf{Boocock 2025 (single-cell eQTL mapping)}\,$\bullet$\newline {\scriptsize Expression / single cell; single-cell transcriptome + eQTL} & P, 1 & 393 segregants, 27,744 cells & glucose / galactose & $5.5\times 10^{4}$$\dagger$ & $3.3\times 10^{8}$ & segregant-\allowbreak WGS & -- & 27,744 single cells drawn from 393 sequenced recombinants in two carbon sources, a median of 17 cells per segregant. High-dimensional on both axes at once, which almost nothing in yeast is, and the closest existing template for the proposed design. \\
\addlinespace[5pt]
2 & \textbf{Hale 2024 (CRISPRi x natural variation)}\,$\bullet$\newline {\scriptsize Natural variation; relative fitness (double-barcode seq)} & P, 1 & 169 segregants x 1,721 genes & 2 induced & $1.4\times 10^{6}$$\dagger$ & $1.4\times 10^{6}$ & segregant-\allowbreak WGS & -- & Directly measures how a genetic perturbation's effect changes with genetic background. The single best dataset for asking whether a model trained on one background transfers to another. \\
\addlinespace[5pt]
3 & \textbf{Jackson 2020 (TF-deletion single-cell atlas)}\,$\bullet$\newline {\scriptsize Expression / single cell; single-cell transcriptome} & P, 1 & 12 genotypes (wild type + 11 TF deletions), 72 barcoded strains & 11 growth conditions & 132$\dagger$ & $7.9\times 10^{5}$ & S288C-\allowbreak KO & -- & The closest thing to a yeast Perturb-seq that exists. A transcribed barcode in the 3' UTR of the marker cassette labels every cell with its genotype, so 72 strains, 12 genotypes at 6 independently constructed replicates, are pooled and read out per cell; every strain appears in every one of the 11 conditions. The perturbations are homozygous diploid whole-ORF deletions in a prototrophic FY4/FY5 background rather than the Giaever collection, which is what permits the minimal and nitrogen-limited media. The DNA-only barcodes of the Giaever collection are the stated reason it could not be used, and that is the design constraint any campaign here inherits. \\
\addlinespace[5pt]
4 & \textbf{Puddu 2019 (WGS of the deletion collection)}\newline {\scriptsize Modality / backbone; structural variants, CNV, aneuploidy} & P, 1 & ~4,800 deletion strains & 1 & $4.8\times 10^{3}$ & $4.8\times 10^{3}$ & S288C-\allowbreak KO & -- & Whole-genome sequence for every strain in the deletion collection. This is what turns the S288C-KO sequence basis from an assumption into a measurement, and it applies to roughly half the rows in this table at once. \\
\addlinespace[5pt]
5 & \textbf{N'Guessan 2025 (segregant scRNA-seq eQTL)}\,$\bullet$\newline {\scriptsize Expression / single cell; single-cell transcriptome + eQTL} & P, 2 & ~4,500 profiled segregants & 1 & $4.5\times 10^{3}$ & $2.7\times 10^{7}$ & segregant-\allowbreak WGS & -- & Single-cell transcriptomes indexed by recombinant natural genotype rather than by an engineered knockout. The closest existing analog to the proposed Perturb-seq design, and the only one that already pairs per-cell expression with sequenced genotypes at scale. \\
\addlinespace[5pt]
6 & \textbf{Hackett 2020 (IDEA inducible-TF transcriptome time series)}\,$\bullet$\newline {\scriptsize Expression / single cell; bulk transcriptome per time point} & P, 2 & over 200 inducible-TF strains & estradiol induction, about eight time points & $1.6\times 10^{3}$$\dagger$ & $9.6\times 10^{6}$ & S288C+\allowbreak designed-\allowbreak edit & about eight time points after induction; one record per strain and time point & Each TF is individually induced from a designed synthetic cassette and the whole transcriptome is followed over time, which is the only yeast dataset where a single-gene perturbation is crossed with a dense time axis. scYeast used it as the target for its time-series perturbation-response task (Fan 2027, Sec. 4.1.7). The instance count treats each strain-by-time-point as a record; per strain it is about 200. \\
\addlinespace[5pt]
7 & \textbf{Hu 2007 (TF deletion expression compendium)}\,$\bullet$\newline {\scriptsize Expression / single cell; mRNA abundance} & P, 2 & 269 TF deletion strains & 1 & 269 & $1.6\times 10^{6}$ & S288C-\allowbreak KO & -- & Transcriptomes for the regulators specifically, predating and complementing the already-built Kemmeren compendium. Deleting a transcription factor is the perturbation whose transcriptome response is most interpretable. \\
\addlinespace[5pt]
8 & \textbf{Dong 2021 (MAGIC + SAM biosensor)}\newline {\scriptsize CRISPR library screen; SAM biosensor fluorescence (FACS)} & P, 2 & ~100,000 guides (a/i/d) & 1 & $1.0\times 10^{5}$$\ddagger$ & $1.0\times 10^{5}$ & S288C+\allowbreak guide & -- & MAGIC re-run with a metabolite biosensor instead of a growth selection, so the label is product concentration rather than fitness. The pattern to copy for every precursor we care about, and the sibling of the already-built Lian 2019 furfural screen. \\
\addlinespace[5pt]
9 & \textbf{Momen-Roknabadi 2020 (inducible CRISPRi library)}\newline {\scriptsize CRISPR library screen; per-guide fitness under induction} & P, 2 & genome-wide guide library & nutrient limitations & $9.0\times 10^{4}$$\ddagger$ & $9.0\times 10^{4}$ & S288C+\allowbreak guide & -- & A second CRISPRi library on a different vector and guide-design rule. Whether a model trained on one library transfers to the other is a cheap, decisive test of guide-level overfitting. Not a Perturb-seq row: one edit per cell and a scalar fitness readout is low-dimensional on both axes, however large the library. \\
\addlinespace[5pt]
10 & \textbf{McGlincy 2021 (genome-scale CRISPRi library)}\,$\bullet$\newline {\scriptsize CRISPR library screen; per-guide fitness} & P, 2 & 61,094 guides, ~10/gene & continuous culture & $6.1\times 10^{4}$ & $6.1\times 10^{4}$ & S288C+\allowbreak guide & -- & Ten guides per gene gives a graded knockdown axis instead of a binary null, and it covers essential genes. This is the library a yeast Perturb-seq would most plausibly be built on. \\
\addlinespace[5pt]
11 & \textbf{Bao 2018 (CHAnGE single-nucleotide library)}\newline {\scriptsize CRISPR library screen; variant fitness (pooled barcode seq)} & P, 2 & tens of thousands of designed variants & inhibitor selection & $6.0\times 10^{4}$$\ddagger$ & $6.0\times 10^{4}$ & S288C+\allowbreak designed-\allowbreak edit & -- & Single-nucleotide resolution rather than whole-gene nulls. The only genotype axis in this table finer than an ORF, and the one that matches how a real strain-design edit is specified. \\
\addlinespace[5pt]
12 & \textbf{Crook 2016 (tunable RNAi, isobutanol + 1-butanol)}\newline {\scriptsize Tolerance / robustness; dose-graded knockdown fitness under alcohol stress} & P, 2 & genome-scale tunable RNAi & isobutanol, 1-butanol & $2.0\times 10^{4}$$\ddagger$ & $2.0\times 10^{4}$ & S288C+\allowbreak guide & -- & Two related alcohols at graded knockdown strength. The dosage axis is the signal a binary screen throws away, and Hsp70 emerged only because of it. \\
\addlinespace[5pt]
13 & \textbf{Roy 2018 (multiplexed precision editing)}\newline {\scriptsize CRISPR library screen; variant fitness} & P, 2 & thousands of designed edits & 1 & $1.6\times 10^{4}$$\ddagger$ & $1.6\times 10^{4}$ & S288C+\allowbreak designed-\allowbreak edit & -- & Genomic rather than plasmid barcodes, which removes the barcode-swapping artifact that quietly corrupts pooled fitness data. The methodological complement to CHAnGE. \\
\addlinespace[5pt]
14 & \textbf{Newman 2006 (protein noise, GFP library)}\,$\bullet$\newline {\scriptsize Modality / backbone; protein abundance and cell-to-cell noise} & P, 2 & ~2,500 GFP-tagged strains & rich / minimal & $5.0\times 10^{3}$$\dagger$ & $1.0\times 10^{4}$ & S288C+\allowbreak tag & -- & Measures the distribution rather than the mean, per protein, in two media. Noise is the quantity a per-cell readout recovers and a bulk one destroys, so this sets what a single-cell design should expect to see. \\
\addlinespace[5pt]
15 & \textbf{Mukherjee 2021 (CRISPRi essential genes x acetic acid)}\newline {\scriptsize Tolerance / robustness; fitness under acid stress} & P, 2 & ~1,100 essential genes & acetic acid vs control & $2.2\times 10^{3}$$\dagger$ & $2.2\times 10^{3}$ & S288C+\allowbreak guide & -- & Acetic acid is the dominant hydrolysate stressor, and this reaches the essential genes a deletion collection structurally cannot. \\
\addlinespace[5pt]
16 & \textbf{Lian 2017 (CRISPR-AID, beta-carotene)}\newline {\scriptsize CRISPR library screen; beta-carotene titer} & P, 2 & combinatorial a/i/d triples & 1 & 200$\ddagger$ & 200 & engineered-\allowbreak chassis & -- & The tri-functional chassis MAGIC was built on, read out on beta-carotene. Small, but it is a combinatorial genotype with a measured product titer, which is the exact record type inverse strain design needs. \\
\addlinespace[5pt]
17 & \textbf{Hughes 2000 (compendium of expression profiles)}\newline {\scriptsize Expression / single cell; mRNA log ratio (two-color array)} & P, 3 & 300 mutations and chemical treatments & 1 & 300 & $1.8\times 10^{6}$ & S288C-\allowbreak KO & -- & The original deletion expression compendium, and the oldest independent measurement of the phenotype the built Kemmeren 2014 measures. Ingesting it makes cross-laboratory replication of the deletion transcriptome testable on gene-by-gene profiles fourteen years apart and on different array platforms, which is a check the built set cannot run against itself. The 300 profiles mix deletions, titratable alleles and compound treatments, so the compound arm carries a reference genotype and the mutant arm an S288C-KO one; the split was not confirmed this pass. \\
\addlinespace[5pt]
18 & \textbf{Lenstra 2011 (chromatin regulator deletion expression)}\,$\bullet$\newline {\scriptsize Expression / single cell; mRNA abundance} & P, 3 & 165 chromatin-regulator deletions & 1 & 165 & $9.9\times 10^{5}$ & S288C-\allowbreak KO & -- & Transcriptomes for the chromatin machinery specifically, from the group behind the built Kemmeren compendium and in the same format. Chromatin regulators are the class whose deletion most reshapes the expression landscape. \\
\addlinespace[5pt]
19 & \textbf{Sun 2013 (mRNA synthesis and decay rates across deletion strains)}\newline {\scriptsize Expression / single cell; mRNA level, synthesis rate and decay rate} & P, 3 & 46 deletion strains of mRNA degradation and metabolism genes & 1 & 46 & $8.3\times 10^{5}$ & S288C-\allowbreak KO & rates derived from metabolic labeling; the rates are stored, not the series & Comparative dynamic transcriptome analysis decomposes a transcript level into the synthesis rate and the decay rate that produce it, across 46 single deletions. That is the transcript-side counterpart of Martin-Perez 2017's protein half-lives, so the two together give both turnover terms. Its finding also warns what a steady-state compendium hides: a change in degradation rate is generally compensated by a change in synthesis rate, so mRNA level is buffered and two strains with the same level can have different kinetics. Overlap with the other deletion-strain expression sets is on the strain axis, not the measurement: the built Kemmeren 2014 (1,484 deletions), Hughes 2000, Hu 2007 and Lenstra 2011 all store a steady-state level per gene, and none stores a rate, so this row is a second layer on shared strains rather than a second copy of one. The 46 strains are single deletions of mRNA degradation and metabolism genes and are expected to be inside the Kemmeren set by gene identity; the exact shared-strain count needs the GEO strain list, which was not fetched, and the abstract does not say the strains came from that collection. \\
\addlinespace[5pt]
20 & \textbf{Jariani 2020 (yeast scRNA-seq through lag phase)}\,$\bullet$\newline {\scriptsize Expression / single cell; single-cell transcriptome} & P, 3 & isogenic wild type & glucose to maltose shift, time course through lag phase & 6$\ddagger$ & $3.6\times 10^{4}$ & reference-\allowbreak only & time course through the lag phase of a glucose-to-maltose shift & A yeast-adapted 10x protocol that resolves the transcriptional states of individual cells during a carbon-source shift, so the readout is the distribution across a population that is not synchronized rather than its mean. One genotype, so the value is the protocol and the per-cell variance structure a campaign has to budget against. \\
\addlinespace[5pt]
21 & \textbf{Guo 2018 (CRISPR-Cas9 tiling of essential genes)}\,$\bullet$\newline {\scriptsize CRISPR library screen; per-guide fitness} & P, 3 & ~3e4 tiling guides & 1 & $3.0\times 10^{4}$$\ddagger$ & $3.0\times 10^{4}$ & S288C+\allowbreak guide & -- & Tiles guides across genes rather than sampling a few per gene, so the readout resolves within-gene functional regions. A finer input axis than a gene-level library. \\
\addlinespace[5pt]
22 & \textbf{Urbonaite 2021 (yeastDrop-Seq under drug treatment)}\,$\bullet$\newline {\scriptsize Expression / single cell; single-cell transcriptome} & P, 3 & isogenic wild type & mycophenolic acid, guanine, both, untreated & 4 & $2.4\times 10^{4}$ & reference-\allowbreak only & -- & A yeast-optimized Drop-seq with a small chemical condition axis, including a two-compound arm. One genotype, so it enters as a platform reference: the pair of single-drug arms against their combination is a per-cell readout of a two-factor environment, which is the environmental analog of the combinatorial genetic perturbation a Perturb-seq needs. \\
\addlinespace[5pt]
23 & \textbf{Nadal-Ribelles 2019 (sensitive yeast scRNA-seq)}\,$\bullet$\newline {\scriptsize Expression / single cell; single-cell transcriptome} & P, 3 & clonal wild type & baseline + stress & 2$\ddagger$ & $1.2\times 10^{4}$ & reference-\allowbreak only & -- & Establishes within-clone transcript correlation, which is the noise floor any per-strain Perturb-seq mean has to clear. Directly sizes the cells-per-guide question. \\
\addlinespace[5pt]
24 & \textbf{Gasch 2017 (single-cell stress heterogeneity)}\,$\bullet$\newline {\scriptsize Expression / single cell; single-cell transcriptome} & P, 3 & wild type & stress vs unstressed & 2$\ddagger$ & $1.2\times 10^{4}$ & reference-\allowbreak only & -- & Separates intrinsic from extrinsic expression heterogeneity under stress. Sets what fraction of Perturb-seq variance is recoverable signal rather than cell-state noise. \\
\addlinespace[5pt]
25 & \textbf{Jaffe 2019 (multiplexed CRISPR interference epistasis)}\,$\bullet$\newline {\scriptsize CRISPR library screen; fitness / genetic interaction} & P, 3 & ~1e4 guide pairs & 1 & $1.0\times 10^{4}$$\ddagger$ & $1.0\times 10^{4}$ & S288C+\allowbreak guide & -- & Two guides per cell, so the perturbation space is combinatorial rather than one edit at a time. This is the shape a Perturb-seq input axis should have, measured here against fitness only. \\
\addlinespace[5pt]
26 & \textbf{Jackson 2023 (wild-type scRNA-seq time course for RNA kinetics)}\,$\bullet$\newline {\scriptsize Expression / single cell; single-cell transcriptome (173,361 cells)} & P, 3 & wild type & one time course, sampled continuously & 1 & $6.0\times 10^{3}$ & reference-\allowbreak only & one time course, sampled continuously & The second scYeast pre-training set: 173,361 cells from one wild-type strain sampled continuously through a time course without metabolic labeling, so the strain-by-condition count is one and the value is the per-cell distribution of transcriptional states, not a genotype axis. A reference backbone for any single-cell expression head, and the companion of Jackson 2020 on the same platform. \\
\addlinespace[5pt]
27 & \textbf{Brettner 2024 (ultra-high-throughput yeast scRNA-seq)}\,$\bullet$\newline {\scriptsize Expression / single cell; single-cell transcriptome} & P, 3 & method & method & --$\ddagger$ & -- & reference-\allowbreak only & -- & The platform the costing model in experiments/024 is built on. Not a dataset to ingest; the reason it is here is that a yeast Perturb-seq proposal has to name its chemistry. \\
\addlinespace[5pt]
28 & \textbf{Mulleder 2012 (prototrophic deletion collection)}\newline {\scriptsize Modality / backbone; strain resource (no phenotype)} & P, 3 & 5,185 prototrophic + 839 titratable essential & 1 & --$\ddagger$ & -- & S288C-\allowbreak KO & -- & The background the already-built Mulleder 2016 and Messner 2023 were measured in, and mulleder2016.py already cites it. Listed not as an ingestion target but because that loader records the prototrophy-restoring markers as unmodeled: they are a GeneAddition this schema does not yet carry, and auxotrophy markers distort metabolite pools. \\
\addlinespace[5pt]
29 & \textbf{Su 2023 (single-cell transcriptomes under four stresses)}\newline {\scriptsize Expression / single cell; single-cell transcriptome (117 cells, full length)} & P, 4 & wild type & isosmotic, hypoosmotic, glucose starvation, amino acid starvation & 4 & $2.4\times 10^{4}$ & reference-\allowbreak only & -- & Full-length single-cell RNA-seq of 117 cells across four stress conditions in one strain; scYeast's stress-state classification benchmark (Fan 2027, Sec. 4.1.6). Tiny, but it is a full-length protocol rather than 3' tag counting, which matters for isoform-level heads. \\
\addlinespace[5pt]
30 & \textbf{Wang 2022 (single-cell transcriptomes across replicative aging)}\newline {\scriptsize Expression / single cell; single-cell transcriptome (125 cells)} & P, 4 & wild type & 2 h, 16 h, 36 h of aging & 3 & $6.6\times 10^{3}$ & reference-\allowbreak only & three ages: 2 h, 16 h and 36 h & Three age groups of one strain, 37, 43 and 45 cells each, about 2,200 genes detected per cell. Small, but it is the only yeast single-cell set with age as the condition axis, and scYeast used it to infer age-specific regulatory relationships (Fan 2027, Sec. 4.1.4). \\
\addlinespace[5pt]
31 & \textbf{Albert 2018 (eQTL in 1,012 segregants)}\,$\bullet$\newline {\scriptsize Natural variation; mRNA abundance (RNA-seq)} & M, 1 & 1,012 sequenced segregants & 1 & $1.0\times 10^{3}$ & $5.8\times 10^{6}$ & segregant-\allowbreak WGS & -- & Transcriptomes on a thousand sequenced recombinants, so genotype maps to a genome-wide expression vector rather than one trait. The single best existing analog to a Perturb-seq readout on a high-dimensional genotype axis. \\
\addlinespace[5pt]
32 & \textbf{Jakobson 2025 (genome-to-proteome map)}\newline {\scriptsize Natural variation; protein abundance, >6,400 pQTL} & M, 1 & 800 F6 segregants & 1 & 800 & $1.6\times 10^{6}$ & segregant-\allowbreak WGS & -- & Natural alleles mapped to enzyme abundance. A lower-risk engineering lever than a heterologous construct, and the proteome counterpart to the metabolite-QTL panels. \\
\addlinespace[5pt]
33 & \textbf{Muenzner 2024 (natural-isolate proteome)}\newline {\scriptsize Natural variation; protein abundance (DIA-MS)} & M, 1 & 796 natural isolates & 1 & 796 & $1.6\times 10^{6}$ & isolate-\allowbreak WGS & -- & Proteomes on the sequenced isolate panel whose transcriptomes Caudal 2024 already supplies. That overlap is the paired anchor the RNA-to-protein inference thesis needs. \\
\addlinespace[5pt]
34 & \textbf{Cooper 2010 (CE-MS amino-acid metabolome)}\newline {\scriptsize Metabolite / precursor; free amino-acid pools (capillary electrophoresis)} & M, 1 & ~4,700 YKO & 1 & $4.7\times 10^{3}$ & $9.4\times 10^{4}$ & S288C-\allowbreak KO & -- & An independent platform measuring the same trait class as the already-built Mulleder 2016 on an overlapping genotype axis. Two independent measurements of one phenotype is the cleanest available test of whether a model has learned biology or a batch. \\
\addlinespace[5pt]
35 & \textbf{Aulakh 2025 (genome-scale ionome)}\newline {\scriptsize Metabolite / precursor; intracellular metal-ion content (ICP-MS)} & M, 1 & ~4,800 YKO & 1 & $4.8\times 10^{3}$ & $4.8\times 10^{4}$ & S288C-\allowbreak KO & -- & The only gene-indexed ionome. Metal-cofactor supply gates iron-sulfur and zinc-dependent pathway enzymes, and no built dataset measures it. \\
\addlinespace[5pt]
36 & \textbf{Teyssonniere 2024 (species-wide proteome against transcriptome)}\newline {\scriptsize Natural variation; protein abundance, paired with transcript abundance} & M, 2 & 942 natural isolates & 1 & 942 & $1.9\times 10^{6}$ & isolate-\allowbreak WGS & -- & Quantitative proteomes for 942 sequenced isolates, the same panel whose transcriptomes the built Caudal 2024 supplies, so the strain-level pairing is one to one rather than an overlap. It is also the published answer to what that pairing shows: mRNA and protein correlate weakly at the population gene level, and the variants associated with protein levels overlap those associated with transcript levels by 3 percent. That makes it the calibration for any RNA-to-protein inference rather than only another proteome. DE-DUPLICATION OPEN: it shares authors and a panel with Muenzner 2024 at 796 isolates, and whether the two are independent acquisitions or one re-served is not established here and must be settled before both are built. \\
\addlinespace[5pt]
37 & \textbf{Grossbach 2022 (BY x RM transcriptome, proteome and phosphoproteome)}\newline {\scriptsize Natural variation; transcript, protein and phosphosite abundance} & M, 2 & 112 genomically defined strains (BY x RM panel) & 1 & 112 & $1.1\times 10^{6}$ & segregant-\allowbreak WGS & -- & The only row that measures transcript, protein and protein phosphorylation on one set of sequenced strains, 1,862 proteins and 2,116 phosphopeptides over 988 proteins. The pairing it yields is per strain rather than per gene: for each of the 112 genotypes the mRNA level and the protein level of the same gene are both observed, which is what turns the RNA-to-protein question from a population correlation into a within-strain residual. It also reports that genetic variants changing a protein's phosphorylation are mostly different from those changing its abundance, so the phosphosite layer is not a proxy for the protein layer. \\
\addlinespace[5pt]
38 & \textbf{Leutert 2023 (phosphoproteome x 101 conditions)}\newline {\scriptsize Expression / single cell; phosphosite abundance} & M, 2 & wild type & 101 conditions & 101 & $5.0\times 10^{5}$ & reference-\allowbreak only & -- & Post-translational enzyme control acts faster than transcription and is invisible to every other layer here. Condition axis only. \\
\addlinespace[5pt]
39 & \textbf{Brauer 2008 (growth-rate-controlled chemostat transcriptome)}\newline {\scriptsize Expression / single cell; mRNA abundance} & M, 2 & wild type & 36 chemostats, 6 limitations & 36 & $2.2\times 10^{5}$ & reference-\allowbreak only & -- & Separates growth rate from nutrient identity by design, which no batch-culture dataset can. A large share of any expression response is a growth-rate confound, and this is the reference for removing it. \\
\addlinespace[5pt]
40 & \textbf{Hackett 2016 (SIMMER multi-omic flux)}\newline {\scriptsize Expression / single cell; paired metabolome + proteome + flux} & M, 2 & wild type & 25 chemostat states & 25 & $7.5\times 10^{4}$ & reference-\allowbreak only & -- & The only jointly measured metabolome, proteome and fluxome. No genotype axis, so it is a mechanism prior rather than training data, but it is the only place the three layers are measured on the same cells. \\
\addlinespace[5pt]
41 & \textbf{Zhu 2014 (kinase / phosphatase lipidomics)}\newline {\scriptsize Metabolite / precursor; lipid class and species abundance} & M, 2 & 129 kinase + phosphatase KOs & 1 & 129 & $2.6\times 10^{4}$ & S288C-\allowbreak KO & -- & A gene-indexed lipidome on a regulatory gene set. Complements the already-built da Silveira 2014 lipidomics, which has no kinase axis, and reads out the malonyl-CoA sink directly. \\
\addlinespace[5pt]
42 & \textbf{Gerke 2017 (urea-cycle mQTL)}\newline {\scriptsize Natural variation; untargeted metabolite abundance} & M, 2 & 147 diploid segregants & 1 & 147 & $1.5\times 10^{4}$ & segregant-\allowbreak WGS & -- & A second mQTL panel from a different cross (oak x wine), so cross-background generalization of metabolite QTL can be tested rather than assumed. \\
\addlinespace[5pt]
43 & \textbf{Ambroset 2014 (metabolite QTL)}\newline {\scriptsize Natural variation; 74 metabolite concentrations (LC-MS/MS)} & M, 2 & ~100 segregants & 1 & 100 & $7.4\times 10^{3}$ & segregant-\allowbreak WGS & -- & The founding metabolite-QTL panel. Pairs with Cooper and Mulleder as the natural-variation view of a phenotype space the deletion collection covers by engineering. \\
\addlinespace[5pt]
44 & \textbf{Blank 2005 (13C metabolic flux)}\newline {\scriptsize Metabolite / precursor; intracellular flux distribution (13C-MFA)} & M, 2 & ~200 viable null mutants & glucose minimal & 200$\ddagger$ & $6.0\times 10^{3}$ & S288C-\allowbreak KO & -- & The only real flux measurement tied to single-gene deletions. Everything else in the built set is a concentration or a titer, which is a proxy for flux rather than flux. \\
\addlinespace[5pt]
45 & \textbf{Skelly 2013 (expression variation across isolates)}\newline {\scriptsize Natural variation; transcriptome + proteome} & M, 3 & 22 sequenced isolates & 1 & 22 & $1.8\times 10^{5}$ & isolate-\allowbreak WGS & -- & Transcriptome and proteome on the same isolates, which is the paired anchor the RNA-to-protein inference argument needs. Few strains, but the pairing is measured rather than joined across studies. \\
\addlinespace[5pt]
46 & \textbf{Airoldi 2016 (nitrogen-limited steady-state and dynamic transcriptome)}\newline {\scriptsize Expression / single cell; bulk transcriptome} & M, 3 & wild type, prototrophic & nitrogen sources x steady-state growth rates, plus an upshift time course & 20$\ddagger$ & $1.2\times 10^{5}$ & reference-\allowbreak only & an upshift time course beside the steady states & Chemostat transcriptomes under controlled nitrogen limitation from the same laboratory, and on the same media axis, as Jackson 2020's NLIM-GLN, NLIM-PRO, NLIM-NH4 and NLIM-UREA conditions. That shared axis is what makes the wild-type baseline of the single-cell atlas separable from the deletion effect, and it is the nitrogen counterpart of Brauer 2008's carbon-limited growth-rate series. Counts are from the citation rather than the released series and need confirming. \\
\addlinespace[5pt]
47 & \textbf{McManus 2014 (ribosome profiling, allele-specific)}\newline {\scriptsize Modality / backbone; ribosome occupancy paired with mRNA abundance} & M, 3 & S. cerevisiae, S. paradoxus and their F1 hybrid & 1 & 3 & $3.3\times 10^{4}$ & reference-\allowbreak only & -- & Translation efficiency measured gene by gene alongside the mRNA abundance of the same gene, which is the first of the two terms that make protein level differ from transcript level. Its result is the reason the term is needed: ribosome occupancy is more conserved than transcript abundance, so translation buffers divergence in mRNA rather than tracking it. The S. paradoxus arm is out of species and would be stored separately or dropped; the cerevisiae arm and the hybrid's cerevisiae alleles are in scope. \\
\addlinespace[5pt]
48 & \textbf{Martin-Perez 2017 (protein half-lives)}\newline {\scriptsize Modality / backbone; protein turnover rate} & M, 3 & wild type, exponential growth & 1 & 1 & $3.2\times 10^{3}$ & reference-\allowbreak only & half-lives derived from a labeling time course; the rates are stored & Turnover rates for 3,160 proteins, the second term that makes protein level differ from transcript level. Half-lives span three orders of magnitude and are not normally distributed, so a single global scaling from mRNA to protein cannot be right, and this is the per-gene coefficient that replaces it. The paper also reports that localization, complex membership and connectivity predict turnover better than sequence features do, which names the covariates such a model needs. \\
\addlinespace[5pt]
49 & \textbf{Boer 2010 (metabolome across nutrient limitations)}\newline {\scriptsize Metabolite / precursor; intracellular metabolite concentrations} & M, 3 & wild type & nutrient limitations x growth rates & 25 & $2.5\times 10^{3}$ & reference-\allowbreak only & -- & The metabolite counterpart to Brauer on the same chemostat design, so metabolite pools are separated from growth rate the same way. Covers central-carbon nodes the deletion metabolomics rows measure at low power. \\
\addlinespace[5pt]
50 & \textbf{Tengolics 2024 (domestication metabolome)}\newline {\scriptsize Natural variation; 19 amino acids + 78 other metabolites} & M, 3 & 17 S. cerevisiae populations & 1 & 17 & $1.6\times 10^{3}$ & isolate-\allowbreak WGS & -- & Shows which metabolic traits already moved under domestication, which is a prior on which are tractable to engineer further. Small n. \\
\addlinespace[5pt]
\midrule \multicolumn{10}{@{}l}{\textbf{Ten extra, in rank order: rows 51--60 replace a row above that proves unreachable.}}\\ \midrule
51 & \textbf{Eder 2020 (flux QTL)}\newline {\scriptsize Natural variation; modeled intracellular flux} & M, 3 & segregant panel (n unconfirmed) & wine fermentation & 100$\ddagger$ & 100 & segregant-\allowbreak WGS & -- & Flux as a QTL phenotype, with PDB1 and VID30 validated. The label is modeled rather than measured, which caps how far it can be trusted. \\
\addlinespace[5pt]
52 & \textbf{Foss 2007 (BY x RM segregant proteome)}\newline {\scriptsize Natural variation; protein abundance (label-free MS)} & M, 3 & BY x RM segregants; the count is not stated in the abstract and was not confirmed & 1 & --$\ddagger$ & -- & segregant-\allowbreak WGS & -- & The first segregant proteome, and the independent replicate of Grossbach 2022 on the same cross fifteen years earlier and on a different platform. Its finding is the one this whole group turns on: the loci that influence protein abundance differ from those that influence transcript levels. Ranked below Grossbach because Grossbach measures the transcript layer on the same strains and this does not, and because neither the segregant count nor the proteome depth was confirmed this pass. \\
\addlinespace[5pt]
53 & \textbf{Yu 2021 (proteome and metabolome under nitrogen limitation)}\newline {\scriptsize Metabolite / precursor; proteome + metabolome} & M, 4 & wild type & nitrogen-limited chemostats & 10$\ddagger$ & $2.5\times 10^{4}$ & reference-\allowbreak only & -- & Nitrogen limitation is the standard trigger for redirecting carbon to product, and this measures both layers under it. Condition axis only, and small. \\
\addlinespace[5pt]
54 & \textbf{de Boer 2020 (100M random promoters)}\,$\bullet$\newline {\scriptsize Regulatory DNA; expression driven (YFP sort-seq)} & S, 1 & >100M random promoters & complex + defined & $1.0\times 10^{8}$ & $1.0\times 10^{8}$ & S288C+\allowbreak reporter-\allowbreak locus & -- & The largest sequence-to-phenotype dataset in yeast by orders of magnitude, and every genotype is an exactly known 80mer. The perturbation space is continuous rather than a gene list, which is what a model of regulatory grammar needs and no deletion collection can supply. \\
\addlinespace[5pt]
55 & \textbf{Vaishnav 2022 (regulatory DNA fitness landscape)}\,$\bullet$\newline {\scriptsize Regulatory DNA; expression + inferred fitness} & S, 1 & ~2e7 promoter variants & 1 & $2.0\times 10^{7}$$\ddagger$ & $2.0\times 10^{7}$ & S288C+\allowbreak reporter-\allowbreak locus & -- & Extends de Boer from measurement to a fitness landscape, and tests native and natural-isolate promoter variants against the random library. The bridge between synthetic sequence space and the natural variation the isolate panels carry. \\
\addlinespace[5pt]
56 & \textbf{Lee 2014 (HIP-HOP fitness signatures)}\,\textsuperscript{\textbf{B}}\newline {\scriptsize Tolerance / robustness; chemogenomic fitness score} & S, 1 & ~11,000 (het + hom) & 3,356 compounds & $1.3\times 10^{7}$ & $1.3\times 10^{7}$ & S288C-\allowbreak KO & -- & Largest chemical-genetic matrix in existence, on the exact YKO genotype axis Costanzo and Kuzmin already use. Biggest single instance-count gain available. BLOCKED: WS15 already attempted it and it stands as awaiting author matrices, so the Science SI and the lab portal did not yield per-strain values. Unblocking is an author request, not a loader. \\
\addlinespace[5pt]
57 & \textbf{Nguyen Ba 2022 (barcoded bulk QTL, 100k segregants)}\,$\bullet$\newline {\scriptsize Natural variation; fitness (Bar-seq)} & S, 1 & ~100,000 barcoded segregants & 8 conditions & $8.0\times 10^{5}$$\dagger$ & $8.0\times 10^{5}$ & segregant-\allowbreak WGS & -- & The largest genotyped recombinant panel in yeast, each segregant barcoded and low-coverage sequenced against known parents. Two orders of magnitude more genotypes than Bloom 2013 on the same kind of cross, which is what makes small-effect and epistatic loci detectable. \\
\addlinespace[5pt]
58 & \textbf{Galardini 2019 (four backgrounds x 38 conditions)}\newline {\scriptsize Natural variation; colony-size fitness (S-score)} & S, 1 & 3,786 KOs x 4 backgrounds & 38 conditions & $5.8\times 10^{5}$$\dagger$ & $5.8\times 10^{5}$ & isolate-\allowbreak WGS & -- & The same knockout built independently in four sequenced backgrounds. Quantifies the 18.5 percent of deletion phenotypes that do not transfer, which is the generalization risk the whole substrate is exposed to. \\
\addlinespace[5pt]
59 & \textbf{Parsons 2006 (bioactive-compound profiling)}\newline {\scriptsize Tolerance / robustness; hypersensitivity score} & S, 1 & ~5,000 viable YKO & 82 compounds & $4.1\times 10^{5}$$\dagger$ & $4.1\times 10^{5}$ & S288C-\allowbreak KO & -- & Full-collection strain coverage against a modest compound panel, the complement to the diagnostic-panel screens above. \\
\addlinespace[5pt]
60 & \textbf{Dutta 2026 (barcoded natural-isolate chemical response)}\newline {\scriptsize Natural variation; pooled-barcode fitness} & S, 1 & 520 natural isolates & >600 compounds & $3.1\times 10^{5}$$\dagger$ & $3.1\times 10^{5}$ & isolate-\allowbreak WGS & -- & Natural sequence diversity crossed with a chemogenomic panel, on isolates drawn from the sequenced 1,011 collection. The same strains already carry Caudal 2024 transcriptomes and Muenzner 2024 proteomes, so it is a three-modality join on one genotype axis. \\
\addlinespace[5pt]
\end{longtable}
\endgroup
\end{landscape}
